\documentclass[11pt,letterpaper]{article}
\PassOptionsToPackage{hyphens}{url}
\usepackage[margin=.83in,top=.8in,bottom=.8in,headheight=15pt]{geometry}
\usepackage[T1]{fontenc}
\usepackage{lmodern,amsmath,amssymb,mathtools,bm,booktabs,array,graphicx,tikz}
\usetikzlibrary{arrows.meta,positioning,calc,decorations.pathreplacing,patterns}
\definecolor{ink}{HTML}{17344D}\definecolor{accent}{HTML}{087E8B}\definecolor{pale}{HTML}{EDF5F5}\definecolor{gold}{HTML}{AD6B19}
\providecommand{\dd}{\,\mathrm{d}}\providecommand{\R}{\mathbb{R}}
\providecommand{\E}{\mathbb{E}}\providecommand{\Var}{\operatorname{Var}}
\providecommand{\Cov}{\operatorname{Cov}}\providecommand{\diag}{\operatorname{diag}}
\providecommand{\norm}[1]{\left\lVert#1\right\rVert}
\providecommand{\abs}[1]{\left|#1\right|}
\usepackage{microtype,hyperref,fancyhdr,enumitem,amsthm,sectsty,needspace}
\allsectionsfont{\raggedright}
\hypersetup{colorlinks=true,linkcolor=ink,urlcolor=accent,pdfauthor={Zan Ahmad},pdftitle={Parallel circulation, shunt flow, and oxygen delivery / Review}}
\setlength{\parindent}{0pt}\setlength{\parskip}{6pt}\linespread{1.07}
\setlist{itemsep=3pt,topsep=4pt,leftmargin=*}
\emergencystretch=2em\clubpenalty=10000\widowpenalty=10000
\pagestyle{fancy}\fancyhf{}\fancyhead[L]{\small Mathematical Cardiology}\fancyhead[R]{\small Lecture 5}
\fancyfoot[L]{\footnotesize Zan Ahmad}\fancyfoot[R]{\footnotesize\thepage}\renewcommand{\headrulewidth}{.3pt}
\newtheorem{proposition}{Proposition}\newtheorem{theorem}{Theorem}\newtheorem{lemma}{Lemma}
\theoremstyle{definition}\newtheorem{definition}{Definition}\newtheorem{example}{Example}
\numberwithin{equation}{section}
\begin{document}
\section*{Parallel circulation, shunt flow, and oxygen delivery --- review sheet}
Use this after attempting the worked examples and practice. Explain each claim in words before calculating.

\section*{Key ideas}
\begin{itemize}
\item Topology determines available routes; pressures and constitutive rules determine flow.
\item A signed negative flow reverses the chosen positive arrow.
\item Parallel branch flows add to total flow.
\item Oxygen content, saturation, delivery and consumption are different quantities.
\item Project ideas are due by meeting 6.
\end{itemize}
\section*{Common errors and repairs}
\textbf{Every circuit segment must carry total pump output.} That holds for the ideal serial loop, not each parallel branch.

\textbf{A blue arrow means blood is blue.} Blood is red; diagram colors are conventions and are not required here.

\textbf{Higher content proves greater delivery.} Multiply content by the relevant flow and state the model assumptions.

\section*{Retrieve without your notes}
\begin{itemize}
\item What happens to the split when both resistances are equal?
\item What extra rule distinguishes a valve from a resistor?
\item Can the parallel toy select a Fontan intervention?
\end{itemize}
\section*{Readiness check}
\begin{itemize}
\item Trace the serial normal route and an explicitly simplified parallel single-pump route.
\item Apply a signed pressure-drop rule and conserve flow at a junction.
\item Calculate an oxygen-content mixture and distinguish content from delivery rate.
\item Identify why topology and synthetic parameters alone cannot select a clinical intervention.
\end{itemize}
\end{document}
